\definesubsubsection{Embed Parameters}
The parameters are processed in the order they are given and can effect both the content of the \g{resulting textual component} as well as its \g{style}. A parameter may also affect the processing of parameters after it. Two general types of parameters are defined: flag parameters and value parameters. Flag parameters are single parameters that add or remove an attribute from the \g{resulting textual component}'s \g{style}. Value parameters add an attribute whose value is determined by the parameter following the current one. The following parameter is then skipped over and thus not processed. \\

An \g{implementation} must at least support the parameters specified in this section if permitted by the output \g{document}, but may add additional parameters the implications of which are completely \g{implementation dependant}. If the output \g{document} does not support a particular parameter, or an unknown parameter is given, a \g{warning} is \glink{signalling}{signalled}. \\

\input{5.9.2.1-float.tex}
\definesubsubsubsection{Width}
\defineidentifier{embed-property-width}{width (<pixels ~n+px>|<percent ~n+\%>)}
\definestyle{embed-property-width}{width: size} \\

Causes the embed content's width to be fixed to the specified size. The size can be given in either \inline$pixels$ or \inline$percent$ where \inline$pixels$ will set the width to be the exact amount of pixels given if the document is viewed at its native resolution. \inline$percent$ will scale the width to the given percentage of the width of the \g{document}. If the \g{document} should not have a width, the \inline$percent$ specification does nothing. Unless the \ident{embed-property-height} is also specified, the embed content's aspect ratio must be preserved. \\

\begin{examples}
  \example{[ image assets/markless-logo.png, width 50px ]}{\includegraphics[width=50px]{assets/markless-logo}}
\end{examples}

%%% Local Variables:
%%% mode: latex
%%% TeX-master: "markless"
%%% TeX-engine: luatex
%%% TeX-command-extra-options: "-shell-escape"
%%% End:

\definesubsubsubsection{Height}
\defineidentifier{embed-property-height}{height (<pixels ~n+px>|<percent ~n+\%>)}
\definestyle{embed-property-height}{height: size} \\

Causes the embed content's height to be fixed to the specified size. The size can be given in either \inline$pixels$ or \inline$percent$ where \inline$pixels$ will set the height to be the exact amount of pixels given if the document is viewed at its native resolution. \inline$percent$ will scale the height to the given percentage of the height of the \g{document}. If the \g{document} should not have a height, the \inline$percent$ specification does nothing. Unless the \ident{embed-property-width} is also specified, the embed content's aspect ratio must be preserved. \\

\begin{examples}
  \example{[ image markless-logo.png, width 50px, height 100px ]}{\includegraphics[width=50px,height=100px]{markless-logo}}
\end{examples}

%%% Local Variables:
%%% mode: latex
%%% TeX-master: "markless"
%%% TeX-engine: luatex
%%% TeX-command-extra-options: "-shell-escape"
%%% End:

\input{5.9.2.4-label.tex}
\definesubsubsubsection{Caption}
\defineidentifier{embed-property-caption}{caption <content .*>} 
\definestyle{embed-property-caption}{interaction: image} \\

Causes a \g{textual component} to be output either alongside or within the embed's \g{textual component}. The text may contain any \gpl{inline directive}. The \g{text} held by the \g{content binding} is outputted to this additional \g{resulting textual component}. The text can only contain \gpl{inline directive}.\\

\begin{examples}
  \longexample{[ image assets/markless-logo.png, caption The //Markless// logo image. ]}{\begin{minipage}{.98\linewidth}
      \begin{center}
        \includegraphics[height=2cm]{assets/markless-logo}
      \end{center}
      \vspace{-1cm}
      \captionof{figure}{The \textit{Markless} logo image.}
    \end{minipage}}
\end{examples}

%%% Local Variables:
%%% mode: latex
%%% TeX-master: "markless"
%%% TeX-engine: luatex
%%% TeX-command-extra-options: "-shell-escape"
%%% End:

\definesubsubsubsection{Description}
\defineidentifier{embed-property-description}{description <content .*>} \\

A textual description of the embed's content for use when the embed content cannot be displayed, or when the \g{user} employs a reader or other form of aid system that relies entirely on textual representation.

%%% Local Variables:
%%% mode: latex
%%% TeX-master: "markless"
%%% TeX-engine: luatex
%%% TeX-command-extra-options: "-shell-escape"
%%% End:

%%% Local Variables:
%%% mode: latex
%%% TeX-master: "markless"
%%% TeX-engine: luatex
%%% TeX-command-extra-options: "-shell-escape"
%%% End:
